\par\medskip\noindent
1.4. \underline{Extensiones de un grupo conmutativo libre}

Sea $I$ un objeto de $\underline T$, y pongamos
\begin{equation}\tag{1.4.1}
 P=\mathbb Z[I],
\end{equation}
el $\mathbb Z$-módulo libre generado por $I$ (SGA 4, IV, 11). A toda
extensión $E$ de $P$ por un $\mathbb Z$-módulo dado $G$ (en lo sucesivo
solo se considerarán extensiones conmutativas de grupos), asociemos el
torsor correspondiente bajo $G_P$, y después su restricción a $I$ por
el morfismo canónico
\[
 I\longrightarrow P=\mathbb Z[I].
\]

%%%%%%%%%%  scan idx 166  |  folio 155  %%%%%%%%%%

Se obtiene así un funtor canónico
\begin{equation}\tag{1.4.2}
 \operatorname{EXT}(P,G)\longrightarrow
 \operatorname{TORS}(I,G_I).
\end{equation}
Afirmo que este funtor es una \underline{equivalencia de categorías}.
Equivalentemente, para todo $G$, las aplicaciones
\begin{equation}\tag{1.4.3}
 \operatorname{Hom}_{\mathbb Z}(P,G)\longrightarrow
 H^0(I,G_I)=G(I)
\end{equation}
y
\begin{equation}\tag{1.4.4}
 \operatorname{Ext}^1_{\mathbb Z}(P,G)\longrightarrow H^1(I,G_I)
\end{equation}
son isomorfismos. Más generalmente, si $A$ es un anillo (no
necesariamente conmutativo) de $\underline T$ y $P=A[I]$ es el
$A$-módulo libre generado por $I$, existen isomorfismos canónicos,
funtoriales en el $A$-módulo variable $G$,
\begin{equation}\tag{1.4.5}
 \operatorname{Ext}^i_A(A[I],G)\xrightarrow{\ \sim\ }H^i(I,G_I).
\end{equation}
Para $A=\mathbb Z$ e $i=0,1$, coinciden con (1.4.3) y (1.4.4). Para
definir (1.4.5), recordemos que su primer miembro es el funtor derivado
del funtor
\[
 G\longmapsto\operatorname{Hom}_A(A[I],G),
\]
que, por la definición de $A[I]$, es isomorfo al funtor
$G\longmapsto G(I)=H^0(I,G_I)$. Como $G\longmapsto G_I$ transforma
módulos inyectivos en módulos inyectivos de $\underline T_{/I}$
(SGA 4, V), se sigue la conclusión. Dejamos al lector la comprobación
de la compatibilidad anunciada con (1.4.3) y (1.4.4).

Finalmente, el funtor (1.4.2) es compatible con las imágenes inversas
por un morfismo de topoi
\[
 f:\underline T'\longrightarrow\underline T.
\]
